\section{Methodology}
\label{sec:methodology}

\subsection{Markov Decision Process Formulation}
\label{subsec:mdp}

Portfolio rebalancing is cast as a discrete-time Markov Decision Process (MDP)
$\mathcal{M} = (\mathcal{S}, \mathcal{A}, \mathcal{R}, \mathcal{P}, \gamma)$,
where $\mathcal{S}$ is the state space, $\mathcal{A}$ is the action space,
$\mathcal{R}$ is the reward function, $\mathcal{P}$ is the transition kernel, and
$\gamma \in (0,1)$ is the discount factor.

\textbf{State space.}
At each rebalancing date $t$, the state $s_t \in \mathcal{S} \subset \mathbb{R}^{225}$
is a concatenation of five groups of features:
\begin{enumerate}
  \item \emph{Normalized log returns} $\tilde{r}_{i,t-k}$ for $i = 1,\dots,10$
        and $k = 0,\dots,19$ (lookback $T = 20$ days), standardized to zero mean
        and unit variance over a rolling window: 200 dimensions.
  \item \emph{Realized volatility} $\hat{\sigma}_{i,t}$ estimated over the same
        20-day window for each asset: 10 dimensions.
  \item \emph{Momentum} $m_{i,t} = \sum_{k=1}^{20} r_{i,t-k}$ for each asset:
        10 dimensions.
  \item \emph{Macroeconomic scalars}: the SELIC target (1 dimension) and the
        bank lending spread (1 dimension), each as a rolling 252-day z-score
        (Section~\ref{subsec:data_macro}).
  \item \emph{SELIC regime one-hot encoding}: three binary indicators
        (Rising, Stable, Falling): 3 dimensions.
\end{enumerate}
The total state dimension is $200 + 10 + 10 + 1 + 1 + 3 = 225$.

\textbf{Action space.}
The action $a_t \in \mathcal{A}$ is a portfolio weight vector
$\mathbf{w}_t = (w_{1,t}, \dots, w_{N,t})$ constrained to the $N$-dimensional
unit simplex:
\begin{equation}
  \mathcal{A} = \Delta^N = \left\{ \mathbf{w} \in \mathbb{R}^N : w_i \geq 0,\;
  \sum_{i=1}^N w_i = 1 \right\},
  \label{eq:action_space}
\end{equation}
with $N = 10$.
This formulation enforces long-only positions and full investment, consistent with
the regulatory constraints applicable to many Brazilian fund structures.

\textbf{Reward function.}
The reward at time $t$ combines a tail-risk penalty with a linear expected-return
component:

\begin{equation}
  r_t = -\text{CVaR}_{\alpha}(\mathbf{R}_{p,t}) + \lambda \, \mathbb{E}[\mathbf{R}_{p,t}]
  \label{eq:cvar_reward}
\end{equation}

where $\mathbf{R}_{p,t} = \mathbf{w}_t^\top \mathbf{r}_t$ is the portfolio return,
$\alpha = 0.95$ is the CVaR confidence level, and $\lambda = 0.5$ controls the
trade-off between downside-risk minimization and expected-return maximization.
Negative CVaR values indicate severe drawdowns, so the minus sign in
Equation~\eqref{eq:cvar_reward} converts the risk penalty into a reward to be
maximized.

The CVaR at confidence level $\alpha$ is defined as:

\begin{equation}
  \text{CVaR}_{\alpha}(X) = -\frac{1}{1-\alpha}
  \int_{-\infty}^{\text{VaR}_\alpha(X)} x \, dF_X(x)
  \label{eq:cvar_def}
\end{equation}

where $\text{VaR}_\alpha(X) = \inf\{x : F_X(x) \geq \alpha\}$ is the
Value-at-Risk threshold and $F_X$ is the cumulative distribution function of $X$.
Operationally, CVaR is estimated as the average of the $(1-\alpha)$ worst returns
obtained by applying the candidate weights $\mathbf{w}_t$ to the trailing window
$\mathbf{r}_{t-T:t}$, with $T = 20$ days.

This deserves emphasis, because it is a property of the implementation rather
than of the CVaR objective in general.
The window scored by the reward is the same window of returns that the state
already encodes in its first 200 dimensions.
The reward is consequently a deterministic function of the state and the action:
it introduces no information beyond what the agent observes, and the transition
carries no feedback about the consequence of the action taken.
The Sharpe-based reward used in configurations A1 and A2, by contrast, is the
realized portfolio return on day $t$ under the behavioral action, which is not
contained in the state.
Section~\ref{subsec:disc_cvar} returns to this asymmetry, which we regard as the
central methodological finding of the ablation.

The discount factor is set to $\gamma = 0.99$, reflecting a preference for near-term
rewards while still valuing long-horizon portfolio performance.

\subsection{Conservative Q-Learning}
\label{subsec:cql}

Standard Q-learning applied to offline data suffers from distributional shift:
the Bellman operator may assign high Q-values to actions that are absent from
the offline dataset and that would perform poorly in deployment.
CQL \citep{Kumar2020} addresses this by augmenting the standard SAC loss with a
conservatism penalty that pushes down Q-values for out-of-distribution
actions while pulling up Q-values for actions observed in the behavioral data.
The CQL loss is:

\begin{equation}
  \mathcal{L}_{\text{CQL}}(\theta) = \mathcal{L}_{\text{SAC}}(\theta)
  + \alpha_{\text{CQL}} \cdot \mathbb{E}_{s \sim \mathcal{D}}\!\left[
    \log \sum_{a} \exp Q_\theta(s, a)
    - \mathbb{E}_{a \sim \beta(s)}\!\left[ Q_\theta(s, a) \right]
  \right]
  \label{eq:cql_loss}
\end{equation}

where $\mathcal{L}_{\text{SAC}}(\theta)$ is the standard Soft Actor-Critic
Bellman residual loss, $\mathcal{D}$ is the offline dataset,
$\beta(s)$ is the behavioral policy distribution at state $s$,
and $\alpha_{\text{CQL}} > 0$ is the conservatism coefficient
(distinct from the CVaR confidence level $\alpha$).

The first term inside the brackets,
$\log \sum_a \exp Q_\theta(s, a)$, is the log-sum-exp over the current policy's
action distribution and approximates the maximum Q-value; minimizing it
discourages the agent from assigning high value to actions not grounded in the
behavioral data.
The second term, $\mathbb{E}_{a \sim \beta(s)}[Q_\theta(s, a)]$, rewards
accurate value estimation for actions that were actually taken by the behavioral
policy, preventing the conservatism penalty from depressing Q-values uniformly.

In financial applications, the conservatism property of CQL is particularly
valuable: it prevents the agent from constructing highly concentrated or
speculative portfolios that would only appear optimal by exploiting spurious
correlations in a finite training sample.
This implicit regularization is analogous to the shrinkage estimators commonly
used in classical mean-variance optimization but is learned end-to-end
through the offline RL training procedure.

The CQL agent is implemented using the \texttt{d3rlpy} library
\citep{Seno2022} with a two-layer actor network of hidden dimension 256 and a
corresponding critic.
Portfolio weights are produced by a softmax output activation that
automatically enforces the simplex constraint in Equation~\eqref{eq:action_space}.
The hyperparameters $\alpha_{\text{CQL}}$ and $\lambda$ are selected via
cross-validation on a held-out validation fold within the training set.

\subsection{Walk-Forward Evaluation Protocol}
\label{subsec:walkforward}

Two evaluation protocols are used in this paper, and they are not the same
protocol applied to different strategies.
The distinction governs how the comparison in Section~\ref{sec:results} should
be read, so it is stated explicitly here.

\textbf{Benchmarks: rolling walk-forward.}
The three classical strategies are evaluated with a rolling walk-forward scheme
over the full 2,487-day sample, partitioned into 35 overlapping folds.
Each fold consists of a 252-trading-day estimation window followed by an
immediately subsequent 63-trading-day test window, corresponding to
approximately one year of estimation data and one quarter of test data.
Windows advance in 63-day increments, so successive estimation sets overlap but
the test windows are non-overlapping and contiguous.
Weights are re-estimated at every fold from the trailing estimation window only.
Concatenating the 35 test windows yields a 2,205-day out-of-sample return series
covering 2016--2024.

\textbf{CQL agent: single chronological holdout.}
The offline RL agent is \emph{not} retrained fold by fold.
The transition dataset is split once in calendar time: transitions from
2016-01-13 to 2021-12-30 ($N = 1{,}485$) form the offline training set
$\mathcal{D}$, generated by the MVP behavioral policy, and transitions from
2022-01-03 to 2024-12-27 ($N_{\text{eval}} = 748$) form a strict holdout that is
touched neither during training nor during hyperparameter selection.
The agent is trained once on $\mathcal{D}$ for a fixed number of gradient steps
and then deployed frozen over the holdout; no weight update occurs after the
calendar cutoff.

Both protocols guarantee that estimation precedes evaluation in calendar time,
so neither admits look-ahead bias.
They are not, however, equally favourable to the strategy being measured.
The benchmarks re-estimate their parameters every quarter and therefore adapt to
the evaluation period as it unfolds, whereas the agent's policy is fixed at the
end of 2021 and must generalise across three subsequent years of unseen regimes.
Any comparison between the two is conservative with respect to the agent in this
specific sense, and a per-fold retrained agent is the like-for-like counterpart
that this study does not provide.

All comparisons reported in Section~\ref{sec:results} are computed on the
intersection of the two evaluation windows, that is, on the trading days for
which the agent and the benchmarks both produce a return.
The benchmark series ends on 2024-11-13, because the walk-forward requires a full
63-day test window and the sample does not accommodate a further fold; the
agent's series runs to 2024-12-27.
Restricting to the common window removes 28 trading days from the agent's series
and is what makes the two sets of numbers comparable at all.

\subsection{Benchmark Strategies}
\label{subsec:benchmarks}

Three classical benchmark strategies are evaluated under the rolling
walk-forward protocol described above:

\textbf{Equal-Weight (EW).} Portfolio weights are fixed at $w_i = 1/N = 0.1$
for all assets at every rebalancing date.
This naive diversification rule \citep{DeMiguel2007} requires no estimation
and serves as a baseline that is robust to estimation error.

\textbf{Minimum-Variance Portfolio (MVP).} Weights are obtained by solving
$\min_{\mathbf{w} \in \Delta^N} \mathbf{w}^\top \hat{\Sigma} \mathbf{w}$,
where $\hat{\Sigma}$ is the sample covariance matrix estimated on the trailing
252-day training window.
The MVP is the classical risk-minimization benchmark and is used as the
behavioral policy to generate the offline dataset.

\textbf{Risk Parity (RP).} Weights are set so that each asset contributes equally
to total portfolio variance.
Risk contributions are equalized iteratively using the algorithm of
\citet{Maillard2010}.

All three benchmarks are rebalanced quarterly, at the fold boundaries of the
walk-forward, and transaction costs are excluded from the primary analysis
to ensure comparability across methods.

\subsection{Ablation Study Design}
\label{subsec:ablation}

To decompose the sources of performance of the full agent, we evaluate four
configurations that differ by exactly one component each.
All four are trained offline, on the same transition dataset, for the same
5{,}000 gradient steps, with the same seed; nothing but the named component
changes between consecutive rungs.

\begin{itemize}
  \item \textbf{A1 -- CQL with conservatism disabled, Sharpe reward, no macro
        state.}
        The conservative weight is set to $\alpha_{\text{CQL}} = 0.001$, which
        reduces the CQL objective in Equation~\eqref{eq:cql_loss} to
        approximately its Soft Actor-Critic component: the penalty on
        out-of-distribution actions is present in form but not in magnitude.
        The reward is the realised daily portfolio return of the behavioral
        action, and the state vector excludes the five macro dimensions.
        This is the least constrained agent in the ladder, and it stands in for
        what an offline actor-critic without a conservatism term produces on
        this dataset.

  \item \textbf{A2 -- CQL with conservatism, Sharpe reward, no macro state.}
        Identical to A1 except that the conservative weight is set to the value
        selected by the grid search of Section~\ref{subsec:cql}.
        Comparing A2 with A1 isolates the contribution of the conservatism
        penalty, not of the offline paradigm: both agents are offline.

  \item \textbf{A3 -- CQL with conservatism, CVaR reward, no macro state.}
        Identical to A2 except for the reward, which becomes the CVaR-based
        objective of Equation~\eqref{eq:cvar_reward}.
        Comparing A3 with A2 isolates tail-risk-aware reward shaping.

  \item \textbf{A4 -- CQL with conservatism, CVaR reward, with macro state.}
        The full proposed model, adding the SELIC one-hot encoding, the SELIC
        target level, and the bank lending spread to the state vector.
        Comparing A4 with A3 isolates the incremental contribution of the macro
        variables.
\end{itemize}

The design therefore attributes performance differences to the conservatism
penalty, the reward function, and the macro state representation independently.
It does \emph{not} speak to the online-versus-offline question: no online agent
is trained in this study, and any claim about the offline paradigm would require
one.
